
Locked joint failures have been studied in  \cite{mu2016kinematic}, \cite{xie2021maximizing}, and \cite{lewis1997fault}. Zonggao et al. \cite{mu2016kinematic} worked on kinematic analysis and fault-tolerant trajectory planning of a non-redundant spatial manipulator under single joint failure. They calculated inverse kinematics using an analytical method. Then they combined it with the desired end-effector performance and employed the concept of fault tolerance angle for trajectory planning under failure. Xie et al. \cite{xie2021maximizing} presents an algorithm to compute the trajectory with the maximum probability of completing a set of point to-point tasks after a locked joint failure. It finds shortest distances between the self-motion manifold joint regions of a point-to-point task. Lewis et al. \cite{lewis1997fault} presented a method to ensure a manipulator can still complete a task even after a locked joint failure. It involves analyzing self-motion manifolds and imposing joint constraints to guarantee reachability. The approach allows for entire regions of the workspace to be reachable even after any joint failure. 

In \cite{banerjee2020tale}, \cite{mcguire2019everybody}, \cite{chang2013robot}, and \cite{brooks2016analysis} failure recovery is accommodated by a human operator. Banerjee et al. \cite{banerjee2020tale} focused on ways of providing decision support to human operators to help robots with failure recovery. The findings indicate that operators should receive both diagnosis and action recommendations for optimal effectiveness in failure recovery. McGuire et al. \cite{mcguire2019everybody} developed an adaptive approach to address failure recovery using a reinforcement learning method. In this approach, when a robot encounters a faulty state, it selects a human assistant to help in recovering from the fault. Chang et al. \cite{chang2013robot} suggest an automatic error recovery approach utilizing Petri nets. This approach enables the robot to recover from known errors and learn how to recover from unknown errors by observing a human demonstrator. Brooks et al. \cite{brooks2016analysis} conducted a survey experiment, they discovered that the perceived effectiveness of recovery strategies, as evaluated by humans, is impacted by the type and severity of the failure as well as the context in which it occurs. 

Du et al. \cite{du2017fault} presented a motion planning technique that enables a six-legged robot to cope with multiple broken legs using null space and Jacobian matrices.

Vonásek et al. \cite{vonasek2015online} proposed a system that employs a high-level motion planner which is equipped with a series of motion primitives, generated by locomotion generators. When a failure happens, the motion planner determines whether to use the existing motion primitives to execute a new plan or to adapt them using particle swarm optimization as necessary.  

Jamshidi et al. \cite{jamshidi2019machine} used Machine Learning to enable run-time self-adaptations and integration of information from multiple heterogeneous models in quantitative planning. Their focus is not on faulty joins but rather on quantitative constraints such as battery power.
\ava{todo for myself: connect to methods}

% One study by Lewis and Maciejewski \cite{r2} proposed a fault-tolerant operation method for kinematically redundant manipulators by identifying areas of the robot's workspace that guarantee failure recoverability using self-motion manifolds. They restricted the joints to certain ranges that maximize the local kinematic failure tolerance measure, which keeps the robot in configurations with large self-motion manifolds. Another study by Patarinsky-Robson et al. \cite{r3} presented an analytical method of recovering from locked joint failure of a 6-degree-of-freedom (dof) robot for failure in three different joints. The approach allowed the 6dof arm to complete a positioning task at any orientation.

% Du and Gao \cite{r5} proposed a method for tolerating multiple faults in a six-legged robot by using null space motion through Jacobian-based motion planning. The study showed that the robot can still locomote with several legs broken. Mu et al. \cite{r6} presented a kinematic analysis and fault-tolerant trajectory planning method for a non-redundant spatial manipulator under a single joint failure using an analytical method for inverse kinematics and maximizing the ability of the space manipulator under a single joint failure.

% Banerjee et al. \cite{r1} investigated the best ways to present a human operator with information about a robot's failure to provide recovery instructions. They presented an experiment in which participants were provided with suggestions for action recommendations and feedback on fault diagnosis. The results suggested that feedforward information (i.e., suggesting action recommendations) was the most useful to humans, and feedback (fault diagnosis suggestions) could supplement that information.

% Finally, McGuire et al. \cite{r7} proposed a reinforcement learning method for a robot to select a human assistant to aid in recovering from a fault. The method utilized a task analysis and monitoring to determine which human assistant would be best suited for the given task.

% In summary, the reviewed studies have proposed different approaches to handle various types of robot failures, including fault-tolerant operation, null space motion planning, and fault-tolerant trajectory planning. They have also investigated how to provide human operators with suggestions for recovery instructions and how to utilize human assistance to aid in fault recovery.


Jacobian projection--based constrained motion planning techniques face similar issues, in addition to failing to generalize to dynamic constraints \ap{cite}, which is the core for most nonprehensile motion primitives.